\section{PPO Implementation and Executed Configuration}
%% ============================================================

The two firms use distinct actor--critic parameters and independent Adam optimizers \citep{schulman2017proximal}. Both firms learn simultaneously in the same market, so single-agent PPO convergence arguments do not imply convergence to a game-theoretic equilibrium. Actors receive neither rivals' prices nor their own realized profits as execution inputs, although each firm's realized rewards train its policy during optimization.

The GRU architecture choice follows directly from the observation structure rather than a claim about commercial pricing software: because the public report averages two consecutive latent price states, adds Gaussian innovation noise, and arrives after a one-step queue delay, a single scalar observation cannot determine whether an adverse report reflects an isolated noise draw or a rival's unilateral price cut \citep{green1984noncooperative, calvano2021algorithmic}. A gated recurrent unit \citep{cho2014gru} provides a compact, differentiable finite-memory sequence encoder over the 32-report window while PPO handles continuous price actions in $[1,2.5]$. Each forward pass recomputes the 128-unit GRU from a fresh zero hidden state across the 32-report window, concatenates normalized episode time $t/512$, and applies a linear projection, $\tanh$, and layer normalization. Thus the implementation is a bounded 32-step sliding-window encoder rather than an unbounded episode-long recurrent state. Policy and value heads share this encoder within each firm, with no weight sharing across firms and no centralized critic. We use GRU+PPO strictly as a controlled testbed for auditing history-dependent continuous pricing policies, and no comparison with an MLP or CNN on the same window is claimed here.

During training, the actor samples a latent Gaussian action and applies $\tanh$ before linearly mapping $[-1,1]$ to prices in $[1,2.5]$ (with action-to-price slope $0.75$). Its log-density includes the analytic $\tanh$ Jacobian correction. This Gaussian action distribution governs training exploration and is distinct from the Gaussian observation noise $\varepsilon_t$, the Gaussian calibration likelihood, and the Gaussian calibration priors (none of which involves Gaussian-process regression or Bayesian optimization). All reported diagnostic curves, final evaluations, and counterfactual audits deploy the frozen deterministic policy $\mu(s)$ obtained by applying $\tanh$ to the latent actor mean. Because $\tanh$ is nonlinear, $\mu(s)$ does not generally equal the expected action under the stochastic training policy; all audit claims apply specifically to this deterministic deployment rule.

\begin{table}[t]
\centering
\caption{Configuration recorded in the completed job files.}
\label{tab:config}
\small
\begin{tabular}{@{}ll@{}}
\toprule
Parameter & Executed Value \\
\midrule
Training seeds in each condition & 20261001 through 20261010 \\
Training interactions per run & \num{1000000} joint market transitions \\
Parallel environments and episode length & 2 environments; 512 periods \\
Public history and encoder width & 32 reports; 128 GRU units \\
Full rollout size & \num{2048} periods per environment; \num{4096} transitions \\
Optimization epochs and minibatch & 10; 256 \\
Adam learning rate and epsilon & \num{0.0003}; \num{0.00001} \\
Discount and GAE parameter & 0.99; 0.95 \\
PPO clipping and entropy coefficient & 0.2; 0.01 \\
Value coefficient and gradient norm cap & 0.5; 0.5 \\
Initial actor log std.\ dev. & $-0.5$; clamped to $[-5, 2]$ during use \\
Informative channel & Delay 1; aggregation 2; added noise scale 0.1 \\
Action-independent channel & Same delay; report $\sim \mathcal{N}(1.6, 0.4^2)$ \\
Diagnostics and final evaluation & 100 checkpoints $\times$ 16 episodes; final 64 episodes \\
Probe family and paired repetitions & 2 firms $\times$ 2 times $\times$ 2 cuts; 128 pairs per probe \\
\bottomrule
\end{tabular}
\end{table}

Advantages are standardized within the complete rollout separately for each firm, using the population standard deviation plus $\epsilon$. For likelihood ratio $r_t$ and standardized advantage $\hat{A}_t$, the minimized loss is
\[
L = -\min(r_t \hat{A}_t, \operatorname{clip}(r_t, 1-\varepsilon, 1+\varepsilon)\hat{A}_t) + c_v \cdot \tfrac{1}{2} \mathrm{MSE} - c_e H.
\]
The effective MSE coefficient is $0.25$ because the code multiplies MSE by $0.5$ and then applies value coefficient $0.5$. Entropy is estimated for the bounded $\tanh$ policy, not its base normal. Stored pre-$\tanh$ actions avoid inverse-$\tanh$ instability.

Episode termination has zero bootstrap; nonterminal rollout boundaries use the next-state value. A shortened final rollout reaches exactly one million interactions. Environments reset synchronously. Deployment uses $\tanh$ of the latent mean; stochastic-policy performance is not reported.

%% ============================================================
